\documentclass[11pt]{article}
\usepackage[margin=1in]{geometry}
\usepackage{amsmath,amssymb,amsthm}
\usepackage{mathtools}
\usepackage{hyperref}

\theoremstyle{plain}
\newtheorem{theorem}{Theorem}
\newtheorem{proposition}{Proposition}
\newtheorem{lemma}{Lemma}
\newtheorem{corollary}{Corollary}
\theoremstyle{definition}
\newtheorem{definition}{Definition}
\newtheorem{remark}{Remark}

\newcommand{\C}{\mathbb{C}}
\newcommand{\dt}{\partial_t}
\newcommand{\dz}{\partial_z}
\newcommand{\abs}[1]{\lvert #1\rvert}

\title{Gershgorin-type velocity bounds for higher-order zero flow}
\author{}
\date{\today}

\begin{document}
\maketitle

We bound the zero velocities of $\dt p=\dz^m p$ by a single, geometrically
explicit quantity built only from inter-zero distances. We call these
``Gershgorin-type'' because, exactly as Gershgorin disks localize eigenvalues
by row-sums of off-diagonal magnitudes, here the velocity of a zero is
controlled by a neighbor-distance sum.

\section{The neighbor-sum and the basic estimate}

Recall (companion note \texttt{derivation\_zero\_odes.tex}) that on simple
times
\begin{equation}\label{eq:ode}
  \dot z_j=-m!\,e_{m-1}\!\Bigl(\bigl\{(z_j-z_a)^{-1}\bigr\}_{a\neq j}\Bigr),
\end{equation}
where $e_{r}$ is the degree-$r$ elementary symmetric polynomial.

\begin{definition}\label{def:Sj}
For a simple configuration $\{z_k\}$ set
\[
  S_j:=\sum_{k\neq j}\frac{1}{\abs{z_j-z_k}},\qquad
  x_k:=\frac{1}{\abs{z_j-z_k}}\ \ (k\neq j),\quad\text{so } S_j=\textstyle\sum_k x_k .
\]
\end{definition}

\begin{lemma}[Triangle inequality on the symmetric sum]\label{lem:tri}
\begin{equation}\label{eq:tri}
  \abs{\dot z_j}
  \le m!\!\!\sum_{\substack{A\subseteq\{1,\dots,N\}\setminus\{j\}\\ |A|=m-1}}
        \;\prod_{a\in A}\frac{1}{\abs{z_j-z_a}}
   = m!\,e_{m-1}(x).
\end{equation}
\end{lemma}

\begin{proof}
Apply $\abs{\cdot}$ to \eqref{eq:ode} and the triangle inequality termwise;
$\abs{\prod_a (z_j-z_a)^{-1}}=\prod_a x_a$.
\end{proof}

\section{Reduction to a single power of $S_j$}

\begin{lemma}[Elementary symmetric vs.\ power of the sum]\label{lem:macl}
For nonnegative reals $x_1,\dots,x_{n}$ and $1\le r\le n$,
\begin{equation}\label{eq:macl}
  e_r(x)\;\le\;\frac{1}{r!}\,\Bigl(\sum_k x_k\Bigr)^{r}.
\end{equation}
\end{lemma}

\begin{proof}
Let $p_1=\sum_k x_k$. Expanding $p_1^{\,r}=\bigl(\sum_k x_k\bigr)^r$ as a sum
over \emph{all} ordered $r$-tuples $(k_1,\dots,k_r)$ gives a sum of nonnegative
terms. The subsum over tuples with \emph{distinct} entries equals
$r!\,e_r(x)$ (each unordered $r$-subset appears in $r!$ orders). Since the
remaining terms (tuples with a repeat) are nonnegative,
$r!\,e_r(x)\le p_1^{\,r}$, which is \eqref{eq:macl}.
\end{proof}

\begin{remark}
Inequality \eqref{eq:macl} is the elementary half of the Maclaurin/Newton
inequalities; we only need the crude form, which holds with the explicit
constant $1/r!$ and requires nothing beyond nonnegativity.
\end{remark}

\begin{proposition}[Master bound]\label{prop:master}
On any simple time, for every $j$,
\begin{equation}\label{eq:master}
  \abs{\dot z_j}\;\le\;m!\,e_{m-1}(x)\;\le\;m\,S_j^{\,m-1}.
\end{equation}
\end{proposition}

\begin{proof}
Combine \eqref{eq:tri} with \eqref{eq:macl} at $r=m-1$:
$m!\,e_{m-1}(x)\le m!\cdot\frac{1}{(m-1)!}S_j^{\,m-1}=m\,S_j^{\,m-1}$.
\end{proof}

\section{The motivating cases}

\begin{corollary}[$m=2$]\label{cor:m2}
\[
  \abs{\dot z_j}\le 2\sum_{k\neq j}\frac{1}{\abs{z_j-z_k}}=2\,S_j .
\]
\end{corollary}

\begin{corollary}[$m=4$, IBM / fourth order]\label{cor:m4}
\begin{equation}\label{eq:m4}
  \abs{\dot z_j}
  \;\le\;24\!\!\sum_{\substack{a<b<c\\ a,b,c\neq j}}
        \frac{1}{\abs{z_j-z_a}\,\abs{z_j-z_b}\,\abs{z_j-z_c}}
  \;\le\;4\,S_j^{\,3}.
\end{equation}
\end{corollary}

\begin{proof}
The first inequality is \eqref{eq:tri} with $m=4$ (so $e_3$ over distinct
triples). For the second, \eqref{eq:macl} at $r=3$ gives
$e_3(x)\le \tfrac{1}{6}S_j^3$, hence $24\,e_3(x)\le 4\,S_j^3$.
\end{proof}

\section{Geometric consequences}

\begin{corollary}[Separation controls speed]\label{cor:sep}
Let $r_j=\min_{k\neq j}\abs{z_j-z_k}$ be the nearest-neighbor distance of $z_j$.
Then $S_j\le (N-1)/r_j$ and therefore
\[
  \abs{\dot z_j}\;\le\;m\left(\frac{N-1}{r_j}\right)^{m-1}.
\]
A well-isolated zero (large $r_j$) is slow; only near-collisions produce large
velocities, and the higher order $m=4$ amplifies near-neighbor effects
cubically rather than linearly.
\end{corollary}

\begin{corollary}[A priori confinement window]\label{cor:window}
Suppose along the flow $r_j(t)\ge\rho>0$ for $t\in[0,T]$. Then
$\abs{z_j(t)-z_j(0)}\le t\cdot m\bigl((N-1)/\rho\bigr)^{m-1}$ for $t\le T$;
in particular each zero stays within a disk whose radius grows at most linearly
in $t$ with the stated explicit rate. (This is a conditional statement: it
presumes the separation lower bound, which is exactly what fails at a
collision.)
\end{corollary}

\begin{remark}[Sharpness]
The bound $m!\,e_{m-1}(x)\le m\,S_j^{\,m-1}$ (for $m>2$) is \emph{saturated} in
spread-out regimes, where the inverse-gap mass is shared among many comparable
terms so that $\max_k x_k/S_j\to0$ and the distinct $(m-1)$-tuples dominate; it is
\emph{poor} when one or a few neighbors dominate (if all but one $x_k\to0$ then
$e_{m-1}\to0$ and the slack is maximal). The spread-out, saturated regime is
exactly the spectral edge (see \texttt{edge\_gershgorin.tex},
Prop.\ on saturation). The intermediate quantity $m!\,e_{m-1}(x)$ in
\eqref{eq:master} is the honest one and is what the code reports alongside the
crude $m\,S_j^{m-1}$.
\end{remark}

\end{document}
